\documentclass[sigconf]{acmart}

\settopmatter{printacmref=false}
\renewcommand\footnotetextcopyrightpermission[1]{}
\pagestyle{plain}

\acmConference[HT '26]{ACM Hypertext Conference}{September 2026}{TBD}
\acmBooktitle{Proceedings of the ACM Hypertext Conference (HT '26)}
\acmDOI{TBD}
\acmISBN{TBD}

\title{Extending the Forward-Forward Algorithm:\\
Goodness Functions, Activation Functions, and Local Spatial Goodness}

\author{Nikhilesh Myanapuri}
\affiliation{%
  \institution{Project Team}
  \country{India}}
\authornote{Roll No: 22110162}

\author{Birudugadda Srivibhav}
\affiliation{%
  \institution{Project Team}
  \country{India}}
\authornote{Roll No: 22110050}

\author{Srivathsa Vamsi Chaturvedula}
\affiliation{%
  \institution{Project Team}
  \country{India}}
\authornote{Roll No: 22110260}

\author{Eswar Ganesh Koleti}
\affiliation{%
  \institution{Project Team}
  \country{India}}
\authornote{Roll No: 22110123}

\begin{document}

\begin{abstract}
The Forward-Forward (FF) algorithm is a local-learning alternative to backpropagation in which each layer is trained using two forward passes: a positive pass and a negative pass. This report provides a detailed extension study of FF across three objectives: (1) goodness-function design, (2) activation-function design, and (3) local spatial goodness using locally connected FF layers. We replicate the FF supervised label-overlay setup on MNIST and run controlled sweeps with deterministic seeding. Our results show that goodness design is the strongest lever in this setting: \texttt{pnorm\_3 + relu} reaches 95.54\% test accuracy, outperforming the baseline \texttt{sum\_sq + relu} at 92.63\%. Huber-like goodness also performs strongly at 93.99\%, indicating that robust growth profiles can improve local discrimination. In activation sweeps (with fixed goodness), ReLU remains the most effective option. For local-spatial FF, however, the best observed setup reaches 57.81\% and underperforms the baseline by 34.82 percentage points. We analyze these outcomes in depth, discussing optimization dynamics, objective geometry, robustness properties, and practical failure modes of spatial-local FF in our current compute and tuning regime.
\end{abstract}

\keywords{Forward-Forward algorithm, biologically plausible learning, goodness functions, activation functions, local learning, MNIST}

\maketitle

\section{Introduction}
Backpropagation has enabled modern deep learning, but it remains a poor mechanistic match to biological learning due to weight transport issues, strict symmetry assumptions, and dependence on global backward error pathways. Hinton's Forward-Forward (FF) algorithm revisits local learning by replacing the forward-backward cycle with two forward passes and layer-wise objectives \cite{hinton2022ff}. Instead of transmitting derivatives through all layers, each layer computes a local goodness signal and updates itself to prefer positive over negative data.

The original FF paper establishes feasibility and identifies a set of open questions critical for scaling and robustness: goodness definition, activation nonlinearity, and local objective structures for spatial data. This report focuses exactly on these unresolved choices.

Our three objectives are:
\begin{enumerate}
  \item \textbf{Objective 1: Goodness-function optimization}. We evaluate a broad family of goodness functions, including squared, unsquared, Huber-like, and p-norm variants.
  \item \textbf{Objective 2: Activation-function engineering}. We test 17 activations, including Student-t negative log-density activations proposed as a promising alternative in the FF paper's future-work discussion.
  \item \textbf{Objective 3: Local spatial goodness}. We build locally connected FF layers and evaluate whether local goodness constraints translate into stronger classification performance.
\end{enumerate}

Beyond reporting metrics, we provide a detailed interpretation of why specific choices worked (or failed), with explicit discussion of objective geometry and optimization behavior.

\section{Background: Forward-Forward Algorithm}
\subsection{Positive and Negative Passes}
In FF, every training step uses two forward passes:
\begin{itemize}
  \item \textbf{Positive pass}: input paired with correct supervisory context (in our supervised setup, correct label overlay).
  \item \textbf{Negative pass}: input paired with incorrect supervisory context.
\end{itemize}
Each layer is encouraged to produce high goodness for positive cases and low goodness for negative cases.

\subsection{Canonical Goodness and Local Objective}
For layer activation vector $\mathbf{y}$, the canonical goodness from the paper is:
\begin{equation}
  g(\mathbf{y}) = \sum_j y_j^2.
\end{equation}

Given positive goodness $g^+$ and negative goodness $g^-$ with threshold $\theta$, a logistic local objective can be written as:
\begin{equation}
\mathcal{L} = \log\left(1 + e^{-(g^+ - \theta)}\right) + \log\left(1 + e^{(g^- - \theta)}\right).
\end{equation}

For supervised inference with label overlay, prediction is obtained by label search:
\begin{equation}
\hat{y} = \arg\max_c \sum_\ell g_\ell(x, c),
\end{equation}
where $g_\ell$ is the goodness from layer $\ell$ for input $x$ with candidate label $c$.

\subsection{Normalization Detail from the Paper}
An important implementation note in the FF paper is that when unsquared activity objectives are used, the normalization mode should be aligned with that objective. In practice, this means L1-style normalization for unsquared formulations and L2-style normalization for squared formulations. This detail matters because mismatch between objective and normalization can distort local gradients and weaken positive-negative separability.

\subsection{Why FF Is Interesting}
FF is attractive because it:
\begin{itemize}
  \item avoids global backward propagation,
  \item supports modular layer-wise updates,
  \item naturally accommodates unknown or non-differentiable forward transforms,
  \item connects to biologically plausible local contrastive learning principles.
\end{itemize}

However, FF performance depends strongly on objective design and data construction quality, making controlled objective studies essential.

\section{Methodology}
\subsection{Implementation}
We implemented a reusable FF codebase in PyTorch with:
\begin{itemize}
  \item goodness function registry,
  \item activation function registry,
  \item fully connected FF layers,
  \item locally connected spatial FF layers,
  \item experiment runners for objective-wise sweeps.
\end{itemize}

The codebase includes deterministic seeding utilities, notebook-driven experiment pipelines, and CSV result export for reproducible reporting.

\subsection{Data and Protocol}
\begin{itemize}
  \item Dataset: MNIST.
  \item For Objectives 1 and 2: fully connected FF network with hidden dimensions $(512, 512, 512)$.
  \item For Objective 3: local spatial FF network with patch-wise local blocks (no weight sharing).
  \item Reproducibility: fixed seeds for Python random, NumPy, and PyTorch.
  \item Strong run settings used in this report: 10 epochs, 30,000 train samples, 10,000 test samples.
\end{itemize}

\subsection{Supervised Label Overlay Setup}
We follow the supervised FF setting where labels are injected into the input representation. For each image, the first 10 positions of the flattened input represent a one-hot label code. Positive data uses the true label; negative data uses an incorrect label sampled uniformly from remaining classes.

\subsection{Objective-Wise Experimental Design}
\paragraph{Objective 1}
Activation is fixed to ReLU. Only goodness is changed.

\paragraph{Objective 2}
Goodness is fixed to \texttt{sum\_sq}. Only activation is changed.

\paragraph{Objective 3}
Locally connected layers process spatial patches. We compare multiple local setup variants and track both test accuracy and local goodness margin.

\subsection{Metrics}
We report:
\begin{itemize}
  \item test accuracy,
  \item test error,
  \item training loss,
  \item positive and negative goodness (Objective 3),
  \item goodness margin (Objective 3).
\end{itemize}

\section{Objective 1: Goodness Function Study}
\subsection{Function Families}
The 14 goodness functions tested can be grouped as follows:
\begin{itemize}
  \item \textbf{Squared family}: \texttt{sum\_sq}, \texttt{mean\_sq}, \texttt{neg\_sum\_sq}, \texttt{sqrt\_sum\_sq}, \texttt{log1p\_sum\_sq}.
  \item \textbf{Unsquared family}: \texttt{sum\_abs}, \texttt{mean\_abs}, \texttt{neg\_sum\_abs}, \texttt{sum\_linear}, \texttt{neg\_sum\_linear}.
  \item \textbf{Robust hybrid}: \texttt{sum\_huber\_like}.
  \item \textbf{p-norm family}: \texttt{pnorm\_1\_5}, \texttt{pnorm\_3}.
  \item \textbf{Smooth nonlinear sum}: \texttt{sum\_softplus}.
\end{itemize}

We compared these while fixing activation to ReLU.

\begin{table*}[t]
\caption{Objective 1 full results (sorted by test accuracy).}
\label{tab:obj1}
\centering
\scriptsize
\begin{tabular}{l l r r r}
\toprule
Goodness & Activation & Train Loss & Test Error & Test Accuracy \\
\midrule
pnorm\_3 & relu & 0.506840 & 0.0446 & 0.9554 \\
sum\_huber\_like & relu & 0.915776 & 0.0601 & 0.9399 \\
sum\_sq & relu & 0.557461 & 0.0737 & 0.9263 \\
neg\_sum\_sq & relu & 0.451928 & 0.0741 & 0.9259 \\
log1p\_sum\_sq & relu & 0.648728 & 0.0777 & 0.9223 \\
sqrt\_sum\_sq & relu & 0.673295 & 0.0941 & 0.9059 \\
mean\_sq & relu & 1.157837 & 0.0988 & 0.9012 \\
mean\_abs & relu & 1.306145 & 0.1519 & 0.8481 \\
neg\_sum\_abs & relu & 1.053865 & 0.2035 & 0.7965 \\
neg\_sum\_linear & relu & 1.110900 & 0.2341 & 0.7659 \\
pnorm\_1\_5 & relu & 1.012079 & 0.3024 & 0.6976 \\
sum\_abs & relu & 1.143102 & 0.5134 & 0.4866 \\
sum\_linear & relu & 1.145557 & 0.5487 & 0.4513 \\
sum\_softplus & relu & 352.891398 & 0.9020 & 0.0980 \\
\bottomrule
\end{tabular}
\end{table*}

\textbf{Key finding:} \texttt{pnorm\_3 + relu} outperformed the baseline \texttt{sum\_sq + relu} by 2.91 accuracy points.

\subsection{Detailed Interpretation: Why \texttt{pnorm\_3} Worked}
We provide four complementary explanations:

\paragraph{(1) Steeper energy landscape for strong features}
For nonnegative activations (ReLU), \(\|y\|_3^3\) grows faster than \(\|y\|_2^2\) for larger magnitudes. This increases relative contribution of highly informative neurons and can sharpen decision contrast in label-search inference.

\paragraph{(2) Better positive-negative separation under local training}
FF training is fundamentally a local separation task. If the objective amplifies discriminative peaks, positive and negative goodness distributions can separate more cleanly, reducing ambiguous label rankings.

\paragraph{(3) Reduced influence of medium-strength responses}
Compared to L2, higher-order norms suppress medium activations in relative terms. This can reduce noisy evidence accumulation from weakly activated units.

\paragraph{(4) Compatibility with ReLU sparsity}
ReLU already promotes sparse, nonnegative representations. A steeper goodness function then acts as a confidence amplifier for sparse high-activation patterns.

\subsection{Detailed Interpretation: Why Huber-Like Goodness Worked}
\texttt{sum\_huber\_like} reached 93.99\% and was second best. A likely reason is robust transition behavior:
\begin{itemize}
  \item near zero: roughly quadratic behavior stabilizes gradients,
  \item at larger magnitudes: more linear growth reduces over-penalization and sensitivity to outliers.
\end{itemize}
In FF, this trade-off can improve stability while preserving meaningful contrast between positive and negative passes.

\subsection{Failure Modes in Objective 1}
Functions like \texttt{sum\_softplus} collapsed severely (9.8\% accuracy). This suggests that merely smooth positive goodness is insufficient; objective geometry must preserve class-discriminative contrast under label search.

\section{Objective 2: Activation Function Study}
We fixed goodness to \texttt{sum\_sq} and compared 17 activation functions.

\begin{table*}[t]
\caption{Objective 2 full results (sorted by test accuracy).}
\label{tab:obj2}
\centering
\scriptsize
\begin{tabular}{l l r r r}
\toprule
Activation & Goodness & Train Loss & Test Error & Test Accuracy \\
\midrule
relu & sum\_sq & 0.557461 & 0.0737 & 0.9263 \\
gelu & sum\_sq & 0.455380 & 0.0794 & 0.9206 \\
silu & sum\_sq & 0.534094 & 0.0862 & 0.9138 \\
leaky\_relu & sum\_sq & 0.412336 & 0.0881 & 0.9119 \\
student\_t\_nll\_df5 & sum\_sq & 0.515548 & 0.0891 & 0.9109 \\
student\_t\_nll\_df3 & sum\_sq & 0.529816 & 0.0931 & 0.9069 \\
student\_t\_nll\_df2 & sum\_sq & 0.541876 & 0.0955 & 0.9045 \\
tanhshrink & sum\_sq & 0.661014 & 0.0968 & 0.9032 \\
mish & sum\_sq & 0.574569 & 0.0988 & 0.9012 \\
bent\_identity & sum\_sq & 1.075029 & 0.1526 & 0.8474 \\
softsign & sum\_sq & 1.109407 & 0.1810 & 0.8190 \\
selu & sum\_sq & 1.040589 & 0.2195 & 0.7805 \\
tanh & sum\_sq & 1.114266 & 0.2211 & 0.7789 \\
elu & sum\_sq & 1.114191 & 0.2973 & 0.7027 \\
hardtanh & sum\_sq & 1.210148 & 0.3102 & 0.6898 \\
softplus & sum\_sq & 1.843963 & 0.7705 & 0.2295 \\
squareplus & sum\_sq & 5.145962 & 0.8865 & 0.1135 \\
\bottomrule
\end{tabular}
\end{table*}

\textbf{Key finding:} ReLU remained the strongest activation in our FF setup.

\subsection{Detailed Interpretation of Objective 2}
\paragraph{Why ReLU remained best}
ReLU likely benefited from three factors: non-saturating positive branch, sparse representation support, and optimization simplicity under local objectives.

\paragraph{Why GELU/SiLU were close but lower}
Smooth activations can improve optimization in backprop settings, but in local FF settings they may soften contrastive goodness responses, leading to slightly weaker separation.

\paragraph{Student-t activations}
Student-t negative log-density activations were competitive but not top. This indicates conceptual promise but also sensitivity to parameterization and threshold coupling.

\paragraph{Low-performing activations}
Very smooth or saturating response profiles (e.g., softplus in this setup) may blur positive-negative contrast, weakening label-search ranking.

\section{Objective 3: Local Spatial Goodness}
We evaluated a local-spatial FF model where each local block computes local goodness. We also tracked goodness margin:
\begin{equation}
\text{goodness margin} = \text{train\_pos\_goodness} - \text{train\_neg\_goodness}.
\end{equation}

\subsection{What Local Spatial Goodness Is Doing}
Instead of a single global goodness over all hidden units, local spatial FF computes goodness per local patch/block. The intended effect is denser local supervision and faster representation shaping in image-like domains.

\subsection{Why This Objective Is Hard}
Local block objectives can improve representational locality but also introduce optimization tensions:
\begin{itemize}
  \item local optima may not align with global class objective,
  \item patch hyperparameters (kernel, stride, channels) strongly affect information flow,
  \item no weight sharing can increase parameter count and training instability.
\end{itemize}

\begin{table*}[t]
\caption{Objective 3 local-spatial results.}
\label{tab:obj3}
\centering
\scriptsize
\begin{tabular}{l r r r r r r r}
\toprule
Setup & Epoch & Train Loss & Pos Goodness & Neg Goodness & Margin & Test Accuracy & Test Error \\
\midrule
unsquared\_goodness & 10 & 1.335932 & 1.094110 & 0.950601 & 0.143510 & 0.5781 & 0.4219 \\
student\_t\_activation & 10 & 1.337863 & 2.092228 & 1.931047 & 0.161182 & 0.5159 & 0.4841 \\
finer\_blocks & 10 & 1.363065 & 2.039513 & 1.966984 & 0.072529 & 0.4870 & 0.5130 \\
coarse\_blocks & 10 & 1.352231 & 2.060941 & 1.947516 & 0.113425 & 0.4453 & 0.5547 \\
\bottomrule
\end{tabular}
\end{table*}

\subsection{Baseline Comparison}
Original fully connected baseline \texttt{sum\_sq + relu}: 0.9263 accuracy.
Best local-spatial setup: 0.5781 accuracy.

\textbf{Gap:} $0.9263 - 0.5781 = 0.3482$ (34.82 percentage points), indicating local-spatial FF underperformed strongly in the current configuration.

\subsection{Detailed Analysis of Objective 3 Outcome}
The current local-spatial setup did not match baseline classification quality. A notable observation is that setups with strong goodness margin did not necessarily yield best test accuracy. This indicates margin is useful as a training diagnostic but is not a standalone proxy for final classification performance under all local architectures.

Potential reasons include:
\begin{itemize}
  \item architectural mismatch between local objectives and global label ranking,
  \item insufficient tuning of local receptive field geometry,
  \item optimization schedule inherited from FC setup not ideal for local model,
  \item under-training relative to complexity introduced by no-weight-sharing local blocks.
\end{itemize}

\section{Discussion}
The experiments collectively show that FF behavior is governed by objective geometry as much as by architecture.

\subsection{Cross-Objective Synthesis}
\begin{itemize}
  \item \textbf{Objective 1}: largest gains came from goodness redesign, confirming that local objective shape is a first-order design choice in FF.
  \item \textbf{Objective 2}: activation choice mattered, but less than goodness choice in our regime.
  \item \textbf{Objective 3}: local-goodness structure requires dedicated tuning; direct transfer of FC settings is insufficient.
\end{itemize}

\subsection{Practical Implication}
For practitioners using FF, the strongest immediate return appears to come from objective engineering (goodness design and threshold coupling) before aggressive architecture changes.

\section{Threats to Validity}
\begin{itemize}
  \item \textbf{Internal validity}: although deterministic seeding reduces stochastic drift, single-seed results can still be brittle.
  \item \textbf{Construct validity}: goodness margin is not always equivalent to final classification quality.
  \item \textbf{External validity}: MNIST findings may not transfer directly to more complex datasets without architectural adaptation.
  \item \textbf{Comparative validity}: exact parity with all paper hyperparameters was not enforced; comparisons are directional, not absolute replication claims.
\end{itemize}

\section{Limitations}
\begin{itemize}
  \item Full paper-scale reproduction was not attempted (capacity and schedule differences remain).
  \item Compute limits constrained exhaustive hyperparameter sweeps.
  \item Most reporting uses fixed-seed comparisons; multi-seed confidence intervals are pending.
  \item Local-spatial architecture likely needs targeted tuning (patch geometry, channel count, thresholding, and optimizer settings).
  \item We did not run full factorial interactions between goodness and activation families.
  \item Local objective thresholds were not independently tuned per architecture depth.
\end{itemize}

\section{Future Work}
\begin{itemize}
  \item Multi-seed statistical analysis with confidence intervals and significance testing.
  \item Targeted tuning of local spatial FF (kernel/stride/channel search and schedule redesign).
  \item Hybrid objectives combining robust goodness and adaptive thresholds.
  \item Extension to CIFAR-scale experiments with stronger data augmentation.
  \item Evaluation of temporal decoupling and wake/sleep-style training variants.
\end{itemize}

\section{Conclusion}
This project provides a detailed extension of the Forward-Forward algorithm across objective, activation, and architecture dimensions. The clearest practical outcome is that objective design is decisive: \texttt{pnorm\_3} and Huber-like goodness substantially improved performance relative to the baseline objective. ReLU remained the strongest activation under fixed goodness. The local-spatial variant, while conceptually aligned with FF's locality motivation, underperformed strongly in current form, highlighting that local objective success depends on careful architecture-optimization co-design. Overall, FF remains a promising framework whose performance envelope is highly design-sensitive and still far from fully explored.

\appendix

\section{Appendix A: Objective 1 Top Comparisons}
\begin{itemize}
  \item \texttt{pnorm\_3 + relu}: 95.54\%.
  \item \texttt{sum\_huber\_like + relu}: 93.99\%.
  \item baseline \texttt{sum\_sq + relu}: 92.63\%.
\end{itemize}

\section{Appendix B: Reproducibility Notes}
All experiments were run from notebook pipelines with deterministic seeding for \texttt{random}, \texttt{numpy}, and \texttt{torch}. Result CSVs were exported from each objective notebook and used directly for the tables in this report.

\begin{thebibliography}{9}

\bibitem{hinton2022ff}
Geoffrey Hinton.
\newblock The Forward-Forward Algorithm: Some Preliminary Investigations.
\newblock \emph{arXiv preprint arXiv:2212.13345}, 2022.

\bibitem{lowe2019}
Sindy Lowe, et al.
\newblock Putting an End to End to End: Gradient-Isolated Learning of Representations.
\newblock \emph{NeurIPS}, 2019.

\bibitem{dropout}
Nitish Srivastava, Geoffrey Hinton, Alex Krizhevsky, Ilya Sutskever, and Ruslan Salakhutdinov.
\newblock Dropout: A Simple Way to Prevent Neural Networks from Overfitting.
\newblock \emph{Journal of Machine Learning Research}, 15(56):1929--1958, 2014.

\bibitem{goodfellow2016}
Ian Goodfellow, Yoshua Bengio, and Aaron Courville.
\newblock \emph{Deep Learning}.
\newblock MIT Press, 2016.

\end{thebibliography}

\end{document}
